\section{Ordered Sets}

A set $S$ is called an \textbf{ordered set} if it is equipped with a relation $<$
such that for every $x, y, z \in S$, the following properties hold:
\begin{itemize}
    \item $x < y \land y < z \implies x < z$
    \item Only one of the relations $x < y$, $x = y$, or $x > y$ holds
\end{itemize}
The relation $<$ is called an \textbf{(total) order} on $S$.

\bigskip

\begin{definition} \label{an:def:order-relations}
    The relations $>$, $\leq$, and $\geq$ are defined as
    \begin{itemize}
        \item $x > y \iff y < x$
        \item $x \geq y \iff x > y \lor x = y$
        \item $x \leq y \iff y \geq x$
    \end{itemize}
\end{definition}

\bigskip

\begin{example}
    $\N, \Z, \Q$ are ordered sets with the following relations:
    \begin{itemize}
        \item $\forall n, m \in \N, m < n \iff \exists k \in \N, n = m + k$
        \item $\forall n, m \in \Z, m < n \iff \exists k \in \Z^{+}, n = m + k$
        \item $\forall p, q \in \Q, q < p \iff \exists r \in \Q^{+}, p = q + r$
    \end{itemize}
    where
    \[
    \Q^{+} \coloneqq \setbuilder*{[(p, q)] \in \Q}{p, q \in \Z^{+}} \subseteq \Q
    \]
\end{example}

\bigskip

\begin{definition} \label{an:def:boundedness}
    Let $S$ be an ordered set and $E \subseteq S$.
    \begin{itemize}
        \item $E$ is said to be \textbf{bounded above} if
        \[
        \exists M \in S, \forall x \in E, x \leq M
        \]
        Such an $M$ is called an \textbf{upper bound} of $E$.
        \item $E$ is said to be \textbf{bounded below} if
        \[
        \exists m \in S, \forall x \in E, m \leq x
        \]
        Such an $m$ is called a \textbf{lower bound} of $E$.
    \end{itemize}
\end{definition}

\bigskip

\begin{theorem} \label{an:thm:bounded-subset-of-Z-in-Z-has-extremum}
    Let $E \subseteq \Z$ be nonempty. Then
    \begin{enumerate}[label=\textbf{(\Alph*)}]
        \item $E$ is bounded below $\implies E$ has a least element
        \item $E$ is bounded above $\implies E$ has a greatest element
    \end{enumerate}
\end{theorem}

\begin{proof}
    \begin{enumerate}[label=\textbf{(\Alph*)}]
        \item Let $E$ be bounded below. Then,
        \[
        \begin{array}{c}
            E \text{ is bounded below} \\
            \Downarrow \\
            \exists m \in \Z, \forall x \in E, m \leq x
        \end{array}
        \]
        Define
        \[
        F \coloneqq \setbuilder*{x - m + 1}{x \in E} \subseteq \N
        \]
        Then
        \[
        \begin{array}{c}
            E \neq \vnot \\
            \Downarrow \\
            F \neq \vnot
        \end{array}
        \]
        By \Cref{an:thm:well-ordering-property},
        \[
        \begin{array}{c}
            \exists n \in F, n = \min F \\
            \Downarrow \\
            \forall x \in E, n \leq x - m + 1 \\
            \Downarrow \\
            \forall x \in E, n + m - 1 \leq x
        \end{array}
        \]
        Since $n = \min F$,
        \[
        \begin{array}{c}
            n \in F \\
            \Downarrow \\
            \exists x_{0} \in E, n = x_{0} - m + 1 \\
            \Downarrow \\
            x_{0} = n + m - 1 = \min E
        \end{array}
        \]
        \item Let $E$ be bounded above. Then,
        \[
        \begin{array}{c}
            E \text{ is bounded above} \\
            \Downarrow \\
            \exists M \in \Z, \forall x \in E, x \leq M
        \end{array}
        \]
        Define
        \[
        F \coloneqq \setbuilder*{-x + M + 1}{x \in E} \subseteq \N
        \]
        Then
        \[
        \begin{array}{c}
            E \neq \vnot \\
            \Downarrow \\
            F \neq \vnot
        \end{array}
        \]
        By \Cref{an:thm:well-ordering-property},
        \[
        \begin{array}{c}
            \exists n \in F, n = \min F \\
            \Downarrow \\
            \forall x \in E, n \leq - x + M + 1 \\
            \Downarrow \\
            \forall x \in E, x \leq M - n + 1 \\
        \end{array}
        \]
        Since $n = \min F$,
        \[
        \begin{array}{c}
            n \in F \\
            \Downarrow \\
            \exists x_{0} \in E, n = -x_{0} + M + 1 \\
            \Downarrow \\
            x_{0} = M - n + 1 = \max E
        \end{array}
        \]
    \end{enumerate}
\end{proof}

\bigskip

\begin{definition} \label{an:def:supremum-infimum}
    Let $S$ be an ordered set, and let $E \subseteq S$ be bounded above.
    An element $\alpha \in S$ is called a \textbf{supremum} (or \emph{least upper bound}) of $E$ if
    \begin{itemize}
        \item $\alpha$ is an upper bound of $E$
        \item $\forall \beta \in S, \beta < \alpha \implies \beta$ is not an upper bound of $E$
    \end{itemize}
    A supremum of $E$ is denoted by
    \[
    \alpha = \sup E
    \]
    Similarly, let $E \subseteq S$ be bounded below.
    An element $\alpha \in S$ is called an \textbf{infimum} (or \emph{greatest lower bound}) of $E$ if
    \begin{itemize}
        \item $\alpha$ is a lower bound of $E$
        \item $\forall \beta \in S, \beta > \alpha \implies \beta$ is not a lower bound of $E$
    \end{itemize}
    An infimum of $E$ is denoted by
    \[
    \alpha = \inf E
    \]
\end{definition}
